% Lecture 02: Объявления, управление потоком и ошибки

\section{Лекция 2. Объявления, поток управления и ошибки}
\label{sec:lecture-02}

Первая программа выполняла почти прямую последовательность команд. Реальная
программа должна хранить состояние, принимать решения, повторять действия и
останавливаться при нужном условии. Поэтому поток управления --- не набор
отдельных ключевых слов, а способ записать алгоритм.

\begin{importantbox}[Что нужно уметь после лекции]
Объявлять и инициализировать простые локальные объекты, строить выражения,
писать ветвления и циклы, проверять ввод, объяснять область видимости и
различать compile-time, linker, runtime, logic errors и undefined behavior.
Итоговая практическая цель --- небольшой корректный консольный инструмент.
\end{importantbox}

\subsection{Имя, объект, значение и состояние}

\begin{cppcode}
int attempts{0};
bool running{true};
double threshold{2.5};
char command{'s'};
\end{cppcode}

\begin{definition}
\textbf{Identifier} --- имя, которым исходный код обозначает сущность.
В примере это \texttt{attempts}, \texttt{running}, \texttt{threshold} и
\texttt{command}. Имя не может начинаться с цифры и не должно совпадать с
ключевым словом.
\end{definition}

\begin{definition}
\textbf{Объект} в терминологии C++ --- область хранения с типом и временем
жизни. \textbf{Значение} --- информация, которую выражение вычисляет или
объект хранит. Переменной будем называть именованный объект, значение которого
можно менять.
\end{definition}

Это точнее бытовой фразы «переменная --- коробка». Коробка полезна как первая
модель, но стандарт не обещает, что локальный объект обязательно лежит в CPU
stack: оптимизатор может держать значение в регистре или вообще устранить
отдельное хранение, сохранив наблюдаемое поведение программы.

Тип \texttt{int} хранит целые из диапазона, зависящего от реализации;
стандарт не обещает ему ровно четыре байта. \texttt{double} --- один из
floating-point типов, а не точные вещественные числа. \texttt{char} хранит
символьные/малые целочисленные значения, \texttt{bool} --- \texttt{false} или
\texttt{true}. Этих фундаментальных типов достаточно для текущих алгоритмов.

\subsection{Declaration, definition, initialization, assignment}

\begin{definition}
\textbf{Declaration} сообщает компилятору имя и свойства сущности. Definition
--- объявление, которое вводит саму сущность; для
обычного локального объекта definition предоставляет хранение.
\end{definition}

\begin{cppcode}
int score{10};       // definition и initialization
score = 12;          // assignment уже существующему объекту
extern int total;    // declaration; definition может быть в другом .cpp
\end{cppcode}

Инициализация создаёт начальное значение в момент начала lifetime объекта.
Присваивание заменяет значение уже живого объекта. Поэтому слова
«инициализируем второй раз» для \texttt{score = 12} неверны.

Основные формы, достаточные сейчас:

\begin{cppcode}
int a = 3;     // copy-initialization
int b(3);      // direct-initialization
int c{3};      // direct-list-initialization
int count{};   // value-initialization: для int получаем 0
\end{cppcode}

Фигурная форма хорошо видит некоторые narrowing conversions:

\begin{cppcode}
int truncated = 3.9; // допустимо с потерей дробной части; warning зависит от флагов
// int rejected{3.9}; // ошибка: list-initialization запрещает narrowing
\end{cppcode}

Не нужно превращать выбор скобок в догму. В учебных локальных объектах
\texttt{\{\}} явно показывает инициализацию и защищает от части сужений.

\begin{cppcode}
const int block_count{15};
// block_count = 16; // compile-time error
\end{cppcode}

\texttt{const} означает, что через это имя нельзя менять объект после
инициализации. Это полезный контракт намерения, а не обещание разместить число
в особой физической памяти.

\subsection{Литералы, выражения и преобразования}

Литерал записывает значение прямо в исходнике: \texttt{42}, \texttt{3.5},
\texttt{'q'}, \texttt{true}. Выражение сочетает литералы, имена и операторы и
вычисляет значение:

\begin{cppcode}
int completed{7};
int total{15};
int remaining = total - completed;
bool finished = remaining == 0;
\end{cppcode}

Арифметика использует \texttt{+}, \texttt{-}, \texttt{*}, \texttt{/} и для
целых \texttt{\%}. При \texttt{7 / 2} оба операнда целые, результат равен 3.
В \texttt{7.0 / 2} целое 2 преобразуется к floating-point, результат 3.5.

\begin{cppcode}
int part{1};
int whole{4};
double wrong = part / whole;       // сначала integer 0, затем double 0.0
double right = 1.0 * part / whole; // 0.25
\end{cppcode}

Неявное \texttt{int} $\rightarrow$ \texttt{double} удобно, но некоторые
большие целые представляются неточно. Обратное преобразование отбрасывает
дробную часть; слишком большое floating-point значение для целевого целого
типа может привести к undefined behavior. Сужение лучше делать заметным и
предварительно проверять диапазон. Подробные casts появятся позже.

Операторы \texttt{==}, \texttt{!=}, \texttt{<}, \texttt{<=}, \texttt{>},
\texttt{>=} возвращают \texttt{bool}. Логические \texttt{\&\&},
\texttt{||}, \texttt{!} означают «и», «или», «не»:

\begin{cppcode}
const bool valid_number = number >= 1 && number <= 15;
const bool should_print = show_all || !pair_complete;
\end{cppcode}

\subsubsection{Short-circuit --- не только оптимизация}

Для \texttt{left \&\& right} правый операнд вычисляется только если левый
оказался \texttt{true}. Для \texttt{left || right} --- только если левый
оказался \texttt{false}. Это гарантия языка и способ безопасно поставить
проверку перед опасной операцией:

\begin{cppcode}
const bool ratio_is_large = denominator != 0 && numerator / denominator > 2;
\end{cppcode}

Если знаменатель ноль, деление не выполняется. Переставить операнды нельзя:
integer division by zero имеет undefined behavior.

\subsection{Ветвление if: порядок проверок проектирует алгоритм}

\begin{cppcode}
if (score < 0) {
    std::cout << "invalid\n";
} else if (score < 50) {
    std::cout << "needs work\n";
} else {
    std::cout << "passed\n";
}
\end{cppcode}

Условия проверяются сверху вниз. Выполняется первое подходящее тело, остальные
ветви цепочки пропускаются. Значит, порядок влияет на корректность:
сначала часто ставят ошибки и более узкие случаи, затем общий случай.

Вложенный \texttt{if} нужен, когда второе решение имеет смысл только внутри
первого:

\begin{cppcode}
if (lecture_exists) {
    if (examples_exist) {
        std::cout << "complete pair\n";
    } else {
        std::cout << "examples missing\n";
    }
}
\end{cppcode}

Фигурные скобки синтаксически необязательны для одного statement, но в нашем
стиле ставятся всегда: добавление второй строки тогда не изменит границы ветви.

\begin{mistakebox}[Присваивание вместо сравнения]
\texttt{if (answer = 0)} присваивает ноль и затем преобразует результат к
\texttt{bool}; это корректный синтаксис, но почти наверняка logic error.
GCC с \texttt{-Wall -Wextra -Wpedantic} реально пишет warning вида
\texttt{suggest parentheses around assignment used as truth value}.
Нужное сравнение: \texttt{if (answer == 0)}.
\end{mistakebox}

Warning не обязан останавливать сборку и не доказывает
ошибки. Но его нужно прочитать и устранить или осознанно обосновать.

\subsection{switch: выбор по дискретному значению}

Меню естественно выражается через \texttt{switch}:

\begin{cppcode}
switch (choice) {
case 1:
    std::cout << "show all\n";
    break;
case 2:
    std::cout << "check one\n";
    break;
case 3:
case 4:
    std::cout << "report\n";
    break;
default:
    std::cout << "unknown command\n";
    break;
}
\end{cppcode}

Выражение \texttt{switch} вычисляется один раз, затем выбирается совпавший
\texttt{case}; \texttt{default} обрабатывает отсутствие совпадения.
\texttt{break} выходит из ближайшего \texttt{switch}. Без него выполнение
продолжается в следующий case --- это fallthrough. Иногда оно намеренно, как
для 3 и 4 выше. Если между непустыми case переход нужен специально, C++17
позволяет отметить его \texttt{[[fallthrough]];}.

\begin{implementationbox}
Компилятор может реализовать \texttt{switch} цепочкой сравнений, jump table
или иначе. Стандарт задаёт поведение, но не требует jump table. Выбираем
\texttt{switch} ради ясности дискретного выбора, не по легенде о скорости.
\end{implementationbox}

Для диапазонов и сложных независимых условий удобнее \texttt{if}; case labels
должны быть подходящими constant expressions и не могут означать диапазон.

\subsection{Циклы: повторение с условием остановки}

Перед циклом ответьте на четыре вопроса: какое состояние меняется, когда
продолжаем, какой шаг приближает остановку и что должно оставаться истинным.

\subsubsection{while и validation loop}

\begin{cppcode}
int number{};
while (!(std::cin >> number) || number < 1 || number > 15) {
    std::cout << "Enter an integer from 1 to 15: ";
    std::cin.clear();
    std::cin.ignore(std::numeric_limits<std::streamsize>::max(), '\n');
}
\end{cppcode}

\texttt{while} сначала проверяет условие, поэтому тело может не выполниться ни
разу. Здесь готовый stream API сообщает об ошибке ввода; \texttt{clear}
сбрасывает состояние ошибки, \texttt{ignore} удаляет неподходящий остаток
строки. Это validation loop: выход возможен только после допустимого значения.

\subsubsection{do while}

\begin{cppcode}
int choice{};
do {
    std::cout << "1 - continue, 0 - exit: ";
    std::cin >> choice;
} while (choice != 0);
\end{cppcode}

\texttt{do while} проверяет условие после тела, поэтому тело выполняется хотя
бы один раз. Это бывает естественно для показа меню, но если ввод может
сломаться, всё равно нужна явная обработка состояния потока.

\subsubsection{for и границы}

\begin{cppcode}
int sum{};                              // accumulator
for (int value = 1; value <= 5; ++value) {
    sum += value;
}
\end{cppcode}

Порядок таков: initialization выполняется один раз; condition проверяется
перед каждой итерацией; затем тело; iteration expression; новая проверка.
Переменная \texttt{value} видна в заголовке и теле, но после цикла её scope
заканчивается. \texttt{sum} --- accumulator: хранит результат уже обработанной
части. Переменная, считающая события, называется counter.

\begin{mistakebox}[Off-by-one]
Для значений 1, 2, 3, 4, 5 нужно \texttt{value <= 5}. Условие
\texttt{value < 5} выполнит четыре итерации. Для индексов последовательности,
напротив, обычно естествен полуинтервал \texttt{0 <= index < size}.
Проверяйте первую, последнюю и пустую итерации отдельно.
\end{mistakebox}

Если шаг не меняет состояние к остановке, например забыто \texttt{++value},
может возникнуть бесконечный цикл. Запись \texttt{while (true)} сама по себе
не ошибка: menu loop может завершаться явным \texttt{break}, но путь выхода
должен быть понятен.

\subsubsection{break, continue, sentinel и поиск}

\texttt{break} завершает ближайший цикл. \texttt{continue} пропускает остаток
текущей итерации и переходит к следующей проверке/шагу. Чрезмерное количество
таких переходов усложняет чтение, но для раннего отсева они полезны:

\begin{cppcode}
for (int value = 1; value <= 20; ++value) {
    if (value % 2 != 0) {
        continue;
    }
    if (value > 10) {
        std::cout << "first suitable: " << value << '\n';
        break;
    }
}
\end{cppcode}

Sentinel-controlled loop заканчивает последовательность по специальному
значению, которое не хранит данные:

\begin{cppcode}
int sum{};
int value{};
while (std::cin >> value && value != -1) {
    sum += value;
}
\end{cppcode}

Здесь \texttt{-1} --- sentinel. Short-circuit гарантирует: сравнение
выполняется только после успешного чтения. Так обрабатывается последовательность
неизвестной длины.

\subsubsection{Loop invariant --- человеческое доказательство}

Инвариант --- утверждение, истинное перед каждой итерацией. Для суммирования:
«\texttt{sum} равен сумме уже обработанных значений». Оно истинно до первой
итерации (обработанных значений нет, \texttt{sum == 0}); тело сохраняет его,
добавляя ровно текущее значение; вместе с условием выхода оно объясняет итог.
Это не формальная магия, а компактный способ проверить алгоритм.

\subsection{Block scope, shadowing и lifetime}

Фигурные скобки создают block scope:

\begin{cppcode}
int count{1};
if (count > 0) {
    int message_count{2};
    std::cout << message_count << '\n';
}
// message_count здесь не виден
\end{cppcode}

Scope отвечает, где имя можно использовать в тексте программы. Lifetime
отвечает, когда объект существует при выполнении. Для простого локального
объекта в блоке lifetime начинается после корректной инициализации и
заканчивается при выходе из блока; эти понятия связаны, но не одинаковы.

\begin{cppcode}
int status{1};
if (ready) {
    int status{2}; // shadowing: скрывает внешнее имя внутри блока
    std::cout << status << '\n';
}
\end{cppcode}

Shadowing допустим, но легко заставляет читать не тот \texttt{status}. В
обычном коде лучше выбрать другое имя. Компиляторы могут иметь отдельные
warnings для shadowing; базовые три флага не обязаны включать их все.

\subsection{Карта ошибок по pipeline}

\textbf{Compile-time error:} исходник не удовлетворяет правилам языка, например присваивание объекту \texttt{const}; executable не создаётся; \textbf{Linker error:} translation unit скомпилирован, но linker не нашёл definition объявленной сущности: обычное \texttt{undefined reference}. Это продолжает pipeline лекции~1; \textbf{Warning:} compiler сообщает о подозрительном, но обычно допустимом коде; сборка может продолжиться; \textbf{Runtime failure:} программа запустилась, но не смогла выполнить задачу: например, не найден корень курса или вход неожиданно закончился. Ненулевой exit status позволяет сообщить об этом окружению; \textbf{Logic error:} программа корректно компилируется и выполняется, но даёт неверный ответ, например сумма 1--5 вычисляется циклом \texttt{value < 5}; \textbf{Undefined behavior (UB):} стандарт не накладывает требований на поведение, например при signed integer overflow или integer division by zero. UB не обязано приводить к crash.

Implementation-defined behavior означает, что реализация выбирает один из
допустимых вариантов и документирует выбор (например, размер \texttt{int}).
Unspecified behavior также допускает несколько вариантов, но реализация не
обязана документировать выбор. Это не синонимы UB.

\begin{ubbox}[Почему «у меня работает» недостаточно]
Неправильная программа с UB может напечатать ожидаемое число, упасть или
изменить поведение после оптимизации. Проверять нужно правила операции, а не
только один запуск на одной машине.
\end{ubbox}

\subsection{Diagnostics: compiler и sanitizer решают разные задачи}

Для GCC и Clang используем минимум:

\begin{shellcode}
g++ -std=c++17 -Wall -Wextra -Wpedantic program.cpp -o program
\end{shellcode}

Компилятор обязан диагностировать некоторые нарушения, а warnings помогают
найти дополнительные подозрительные места. Но корректная сборка не доказывает
корректность алгоритма.

UndefinedBehaviorSanitizer добавляет runtime-проверки некоторых классов UB:

\begin{shellcode}
g++ -std=c++17 -fsanitize=undefined \
    -fno-sanitize-recover=undefined program.cpp -o program
./program
\end{shellcode}

Это режим конкретного toolchain, а не часть стандарта C++. Он видит только
проверяемые события реально пройденного пути. Отсутствие diagnostics не
доказывает отсутствие UB во всей программе.

\subsection{Эксперимент 1: course\_status}

Инструмент
\texttt{examples/02\_declarations\_flow\_and\_errors/01\_course\_status}
read-only проверяет все 15 известных пар
\texttt{lectures/NN\_slug.tex} $\leftrightarrow$
\texttt{examples/NN\_slug/}. Он ищет корень курса от текущего каталога вверх,
но ничего не создаёт, не удаляет и не переписывает.

\begin{shellcode}
cmake -S examples/02_declarations_flow_and_errors/01_course_status \
      -B examples/02_declarations_flow_and_errors/01_course_status/build
cmake --build examples/02_declarations_flow_and_errors/01_course_status/build
./examples/02_declarations_flow_and_errors/01_course_status/build/course_status
\end{shellcode}

Программа хранит \texttt{running} и выбранный пункт меню. Внешний
\texttt{while} повторяет меню, \texttt{switch} выбирает действие, а
\texttt{for} проверяет 15 пар. Результаты двух проверок имеют тип
\texttt{bool}; локальные блоки case
ограничивают scope временных имён. \texttt{continue} возвращает к меню после
плохого ввода, \texttt{break} завершает case. Штатный выход имеет status 0,
невозможность найти корень --- 2.

\begin{standardbox}[Язык и библиотека не смешиваем]
\texttt{if}, \texttt{switch}, \texttt{while}, \texttt{for} и scope --- тема
языка этой лекции. \texttt{std::filesystem} --- готовый library API C++17,
которым мы спрашиваем о путях; внутреннее устройство API сейчас не изучаем.
\end{standardbox}

Правильный результат: пункты 1 и 2 сообщают реальное состояние файлов, пункт
3 показывает только неполные пары (для утверждённой структуры их нет), буква
вместо номера не ломает menu loop, 0 завершает программу.

\subsection{Эксперимент 2: control flow и реальные diagnostics}

Второй каталог блока, \texttt{02\_control\_flow\_lab}, содержит несколько
маленьких cases вместо одного искусственного файла.

\begin{shellcode}
printf 'y\n' | \
  examples/02_declarations_flow_and_errors/02_control_flow_lab/run_lab.sh
\end{shellcode}

Лаборатория реально проверяет три ветви \texttt{if}, границу цикла и
short-circuit. Затем она ожидает: ненулевой exit status у демонстрации logic
error; warning для assignment в условии; отказ компиляции при изменении
\texttt{const}; runtime-диагностику UBSan для signed overflow. Expected failure
--- PASS только когда log содержит именно ожидаемую причину. Полные сообщения
конкретного установленного компилятора сохраняются в локальном
\texttt{generated/} и не выдаются за универсальный текст стандарта.

\subsection{Сильные стороны и цена свободы}

C++ прямо выражает control flow, а информация о типах позволяет найти много
ошибок до запуска. Состояние и шаги алгоритма видны в коде; стандартная
библиотека даёт готовые абстракции вроде streams, vector и filesystem. Один
исходник можно собирать в обычном режиме, с усиленными warnings или sanitizer.

Цена этой свободы технически существенна: implicit conversions способны
терять данные; logic errors compiler находить не обязан; некоторые неверные
программы успешно компилируются; язык допускает UB. Поэтому программист
отвечает за условия операций, границы циклов, progress к остановке и проверку
диагностик.

\subsection{Проверка понимания}

После лекции нужно уметь без подсказки:

1) отличить initialization от assignment и declaration от definition; 2) объяснить тип результата сравнения и short-circuit; 3) выбрать \texttt{if} или \texttt{switch} и обосновать порядок проверок; 4) записать validation, sentinel и menu loops; 5) назвать initialization, condition, body и iteration expression цикла; 6) найти off-by-one и сформулировать простой loop invariant; 7) отличить scope от lifetime и заметить shadowing; 8) классифицировать compile, link, runtime, logic errors и UB; 9) объяснить, почему warnings и sanitizer полезны, но не доказывают корректность; 10) написать небольшую консольную программу с корректным exit status.

\subsection{Источники для сверки}

Правила declarations, statements, conversions и undefined behavior сверены с
рабочим draft C++17 N4659. Педагогический порядок условий и циклов сопоставлен
с MIT 6.096 Lecture 2, но старый стиль примеров заменён современным явным
\texttt{std::} и фигурными скобками. Практика warnings и стадий сборки сверена
с \textit{An Introduction to GCC}; target-based сборка продолжает подход
первой лекции и \textit{Modern CMake}.
